\scanpage{5}
$W$ vecteurs de Witt sur corps parfait $k$.

$\nu\geq1$ entier.

On s'intéresse à la catégorie des \deleted{algèbres} \inserted{anneaux} artiniens \deleted{\unsure{sur} algèbres sur} \inserted{de corps rés.} $k$ à idéal maximal muni de puiss. divisées $\nu$-nilpotentes. On va construire un objet initial de cette catégorie.

Pour ceci, on note que les $\Lambda$ sont des algèbres sur $W$ (\unsure{Cohen}). et que l'élément $\lambda_0=p\in W$ s'envoie dans $\unsure{\mathfrak{m}_{\Lambda}|\mathfrak{m}\Lambda}$. Donc $\Lambda$ reçoit
\[
\begin{split}
W^{(\nu)}=W[\Lambda_0,\ldots,\Lambda_\nu]/(&\Lambda_0^p-p!\Lambda_1,\Lambda_1^p-\frac{p^2!}{(p!)^p}\Lambda_2,\ldots,\Lambda_{\nu-1}^p-\frac{p^\nu!}{(p^{\nu-1}!)^p}\Lambda_\nu,\Lambda_\nu^p;\\
&\Lambda_0-\deleted{\unsure{p}}\lambda_0)
\end{split}
\]
[\deleted{\unsure{anneau avec idt}} \inserted{$W$-algèbre} avec idéal $J=(\Lambda_0,\ldots\Lambda_\nu)$ à puiss. div. $\nu$-nilpotentes, \uline{universel} pour les algèbres sur $W$, avec idéal à puissances divisées $\nu$-nilpotentes contenant l'image de $\lambda_0$]. NB. $\Lambda_\nu=\gamma^{p^\nu}(\Lambda_0)$; on utilise le fait que dans $W$ les nbs. premiers $l\ne p$ sont inversibles.

On \deleted{$\unsure{W^{(\nu)}}$} on voit de suite que $W^{(\nu)}$ est artinien de corps résiduel $k$, \deleted{(\unsure{c'est un vrai anneau de} \unsure{mettre|enlever} la relation $\Lambda_0-\lambda_0$), \unsure{qui}} \inserted{l'idéal maximal étant précisément $J$}

On va calculer explicitement $W^{(\nu)}$. Pour ceci, \deleted{\unsure{on pose}} on commence par l'anneau plus simple
\[
\deleted{\unsure{W}}\quad A^{(\nu)}=W[\Lambda_0,\ldots\Lambda_\nu]/(\Lambda_0^p-p!\Lambda_1,\ldots,\Lambda_\nu^p)
\]
(on omet la relation $\Lambda_0-\lambda_0$). On pose aussi
\[
\lambda_r=\frac{p^{p^r}}{p^r!}=\gamma^{p^r}_{\text{habituel}}(p)
\]
et on pose
\[
\Lambda_i=\lambda_i+\Theta_i\qquad 0\leq i\leq\nu.
\]
Par abus de \unsure{langage}, on désignera par la suite par $\Lambda_i$, $\Theta_i$ les images de $\Lambda_i$, $\Theta_i$ dans $A^{(\nu)}$. Alors $A^{(\nu)}$ \unsure{est} une
\mspage{4}
